% GENERATED FILE. Edit canonical lessons, then run npm run generate:books.
% canonical-source-hash: fnv1a64:529c629724c85799
% canonical-lessons: GE-C256-wie-komme-ich, GE-C256-links, GE-C256-review

\chapter{Finding the Way — wie komme ich zum Bahnhof}
\label{ch:ge-finding-the-way}
\hlchaptermodality{256}

\begin{chapteropening}
\textbf{By the end of this chapter:} \emph{I can ask the way with wie komme ich zum or zur, and follow an answer with geradeaus, links and rechts.}

Complete GE-C256-review: ask the way to two places, one with zum and one with zur, and give the answer back in the Sie form.
\end{chapteropening}

\section[wie komme ich zum Bahnhof]{wie komme ich zum Bahnhof — asking the way}

\label{lesson:GE-C256-wie-komme-ich}



\begin{quote}
To get a stranger's attention in the street, say \emph{Entschuldigung!}
\end{quote}

\subsection*{You'll want to know: wie komme ich zum Bahnhof}
\textbf{Entschuldigung, wie komme ich zum Bahnhof?} — \textquotedblleft{}Excuse me, how do I get to the station?\textquotedblright{}

German asks with \emph{kommen}: literally \textquotedblleft{}how do I come to the station?\textquotedblright{} English says \textquotedblleft{}get\textquotedblright{}; German keeps the plain verb of coming.

Two short follow-ups help:

\begin{quote}
\emph{Ist das weit?} — Is it far? \emph{Danke schön!} — Thank you.
\end{quote}

To a stranger, keep to \emph{Sie}: \emph{Können Sie mir helfen?}

\begin{grammarlens}[title={Grammar Lens: zum or zur}]
\textbf{zu}, \textquotedblleft{}to\textquotedblright{}, takes the dative, and it fuses with the article:

\begin{quote}
\textbf{zum} = \emph{zu dem}, for \emph{der} and \emph{das} words: \emph{zum Bahnhof}, \emph{zum Museum} \textbf{zur} = \emph{zu der}, for \emph{die} words: \emph{zur Post}, \emph{zur Apotheke}, \emph{zur Haltestelle}
\end{quote}

So the gender you learned with each noun decides the little word in front of it. \emph{Der Bahnhof}, \emph{zum Bahnhof}. \emph{Die Post}, \emph{zur Post}.

In Switzerland you will hear \emph{Exgüsi} in dialect for \textquotedblleft{}excuse me\textquotedblright{}, usually traced to French \emph{excusez}; \emph{Entschuldigung} is understood everywhere.
\end{grammarlens}

\begin{cousinweb}
\textbf{kommen} $\leftarrow$ Old High German \emph{kweman} $\leftarrow$ Proto-Germanic \emph{*kwemaną}, \textquotedblleft{}to come\textquotedblright{} (EtymDB 2.1): the same verb as English \emph{come}. \textbf{zu} is usually traced to Old High German \emph{zuo}, a cousin of English \emph{to}; EtymDB lists \emph{zu} without a parent, so that is given as the usual account.
\end{cousinweb}

\subsection*{Your turn}
Say these aloud:

\begin{itemize}
\raggedright
  \item \emph{Entschuldigung, wie komme ich zum Bahnhof?}
  \item \emph{Wie komme ich zur Post?} Then \emph{zum Museum}, \emph{zur Haltestelle}
  \item the follow-up: \emph{Ist das weit?}
\end{itemize}

\begin{tcolorbox}[breakable,colback=teal!4,colframe=teal!35!black,title={Before you move on}]
\textquotedblleft{}How do I get to the station?\textquotedblright{} (\emph{\textbf{Wie komme ich zum Bahnhof?}}) \emph{Zum} is short for what? (\emph{\textbf{Zu dem.}}) \textquotedblleft{}To the post office\textquotedblright{}? (\emph{\textbf{Zur Post}}, because it is \emph{die Post}.)
\end{tcolorbox}
\section[Gehen Sie geradeaus, dann links]{Gehen Sie geradeaus — following the answer}

\label{lesson:GE-C256-links}



\begin{quote}
Ask the way to the station. (\emph{Entschuldigung, wie komme ich zum Bahnhof?})

Now the stranger answers, and you have to follow.
\end{quote}

\subsection*{You'll want to know: links}
You have \textbf{geradeaus}, straight on, and \textbf{rechts}, right. The missing one is \textbf{links}, \textquotedblleft{}left\textquotedblright{}: \emph{links}, with a short \emph{i}.

\begin{quote}
\emph{geradeaus} — straight on \emph{rechts} — right \emph{links} — left
\end{quote}

German speakers often add \emph{nach} in front: \emph{nach links}, \emph{nach rechts}, \textquotedblleft{}to the left, to the right\textquotedblright{}. Both ways are standard.

\begin{grammarlens}[title={Grammar Lens: the Sie command, verb first}]
A stranger will tell you the way with \emph{Sie}. The polite command puts the verb first, in its \emph{Sie} form, and \emph{Sie} straight after it:

\begin{quote}
\emph{Gehen Sie geradeaus.} — Go straight on. \emph{Gehen Sie bis zur Ampel, dann links.} — Go as far as the lights, then left. \emph{Nehmen Sie die U-Bahn.} — Take the underground.
\end{quote}

\emph{Bis zur Ampel} is \emph{bis} + \emph{zur}: as far as the lights. \emph{Bis zur Kreuzung}, as far as the crossroads.

The answer may end with the place itself: \emph{Der Bahnhof ist gleich gegenüber}, the station is right opposite.

With a friend, the \emph{du} command you met in chapter forty-six drops the pronoun: \emph{Geh geradeaus!} With a stranger, keep the \emph{Sie}.
\end{grammarlens}

\begin{cousinweb}
\textbf{links} comes from Middle High German \emph{linc}, \textquotedblleft{}left\textquotedblright{}, and EtymDB 2.1 sets Old High German \emph{linka}, \textquotedblleft{}the left hand\textquotedblright{}, beside it; further back its origin is not settled. \textbf{rechts} is built on \emph{recht}, which EtymDB traces through Old High German \emph{reht} to Proto-Germanic \emph{*rehtaz}, \textquotedblleft{}right, direct\textquotedblright{}, the parent of English \emph{right}. The \emph{-s} at the end of both is usually explained as an old adverb ending, the one that also sits in \emph{abends}, \textquotedblleft{}in the evening\textquotedblright{}.
\end{cousinweb}

\subsection*{Your turn}
Say these aloud:

\begin{itemize}
\raggedright
  \item \emph{geradeaus, rechts, links}
  \item \emph{Gehen Sie geradeaus.}
  \item \emph{Gehen Sie bis zur Ampel, dann links.}
  \item \emph{Nehmen Sie die U-Bahn.}
  \item to a friend: \emph{Geh geradeaus, dann rechts!}
\end{itemize}

\begin{tcolorbox}[breakable,colback=teal!4,colframe=teal!35!black,title={Before you move on}]
\textquotedblleft{}Left\textquotedblright{}? (\emph{\textbf{Links.}}) \textquotedblleft{}Go straight on\textquotedblright{}, to a stranger? (\emph{\textbf{Gehen Sie geradeaus.}}) Where does the verb stand in a command? (\textbf{First.}) \textquotedblleft{}As far as the lights\textquotedblright{}? (\emph{\textbf{Bis zur Ampel.}})
\end{tcolorbox}
\section[(review)]{Review — the way there}

\label{lesson:GE-C256-review}



\begin{quote}
Ask the way to the pharmacy. (\emph{Wie komme ich zur Apotheke?}) \emph{Die Apotheke}, so \emph{zur}.
\end{quote}

\begin{grammarlens}[title={What you've built}]
Two lessons, and a whole exchange in the street. You ask with \emph{wie komme ich}, and the gender of the place picks \emph{zum} or \emph{zur}. The answer comes back as a \emph{Sie} command with the verb first, and three short words do much of the work: \emph{geradeaus}, \emph{rechts}, \emph{links}.

One exchange, both halves:

\begin{quote}
\emph{Entschuldigung, wie komme ich zum Museum?} \emph{Gehen Sie geradeaus bis zur Kreuzung, dann rechts. Das Museum ist gleich gegenüber.}
\end{quote}
\end{grammarlens}

\subsection*{Your turn}
Say these aloud:

\begin{itemize}
\raggedright
  \item \emph{Entschuldigung, wie komme ich zum Museum?}
  \item the answer: \emph{Gehen Sie geradeaus bis zur Kreuzung, dann rechts.}
  \item \emph{Wie komme ich zur Haltestelle?}
  \item an answer of your own, with \emph{links}
  \item \emph{Danke schön!}
\end{itemize}

\begin{tcolorbox}[breakable,colback=teal!4,colframe=teal!35!black,title={Before you move on}]
\emph{Der Bahnhof}: \emph{zum} or \emph{zur}? (\emph{\textbf{Zum.}}) \emph{Die Post}? (\emph{\textbf{Zur.}}) \textquotedblleft{}Then left\textquotedblright{}? (\emph{\textbf{Dann links.}})
\end{tcolorbox}
